%% kapa: exam=1 round=23 year=2012 period=1 form=A subject=civil_law subject_ko=민법 seq=11 number=11
%% source: Q-Net 감정평가사 기출문제 (cst00309) artlSeq=5052578 file=raw/1008/5052578/2012년 제23회 감정평가사 제1차시험 기출문제(1교시 A형).pdf
%% conversion: kordoc 4.24.51 -> tools/gen1.py (auto); review=proofread-sonnet; flags=none
\begin{question}{11}
유치권에 관한 설명으로 옳지 않은 것은? (다툼이 있으면 판례에 의함)
\begin{choices}
  \item 유치권자는 피담보채권 전부의 변제를 받을 때까지 경매절차의 매수인에 대하여도 목적물의 인도를 거절할 수 있다.
  \item 유치권자의 점유는 직접점유이든 간접점유이든 관계가 없으나, 그 직접점유자가 채무자인 경우에는 유치권의 요건으로서의 점유에 해당되지 않는다.
  \item 유치권을 행사하고 있는 동안에는 피담보채권의 소멸시효는 진행하지 않는다.
  \item 건물의 임차인은 임대차 존속기간 중 지출한 유익비의 상환을 받을 때까지 유치권을 행사할 수 있다.
  \item 건축업자가 건물을 점유하기 전에 건물공사로 취득한 건축비채권이라도 그 후 그 건물을 점유하였다면 유치권은 성립한다.
\end{choices}
\end{question}
\answer{3}
